\section{Yang--Mills theory and the energy--momentum tensor}
\label{sec:2}
\subsection{The energy--momentum tensor with dimensional regularization}
In the present paper, we consider the $SU(N)$ Yang--Mills theory defined in a
$D$~dimensional Euclidean space. The action is given by
\begin{equation}
   S=\frac{1}{4g_0^2}\int d^Dx\,F_{\mu\nu}^a(x)F_{\mu\nu}^a(x),
\label{eq:(2.1)}
\end{equation}
from the Yang--Mills field strength
\begin{equation}
   F_{\mu\nu}(x)
   =\partial_\mu A_\nu(x)-\partial_\nu A_\mu(x)+[A_\mu(x),A_\nu(x)].
\label{eq:(2.2)}
\end{equation}
We set
\begin{equation}
   D=4-2\epsilon,
\label{eq:(2.3)}
\end{equation}
and then the mass dimension of the bare gauge coupling~$g_0$ is~$\epsilon$.

Assuming that the theory is regularized by the dimensional regularization (for
a very nice exposition, see Ref.~\cite{Collins:1984xc}), one can define the
energy--momentum tensor for the system~\eqref{eq:(2.1)} simply by (see, e.g.,
Ref.~\cite{Freedman:1974gs})
\begin{equation}
   T_{\mu\nu}(x)=\frac{1}{g_0^2}
   \left[F_{\mu\rho}^a(x)F_{\nu\rho}^a(x)
   -\frac{1}{4}\delta_{\mu\nu}F_{\rho\sigma}^a(x)F_{\rho\sigma}^a(x)\right],
\label{eq:(2.4)}
\end{equation}
up to terms attributed to the gauge fixing and the Faddeev--Popov ghost fields,
which are irrelevant in correlation functions of gauge-invariant operators.
Note that the mass dimension of the energy--momentum tensor is~$D$.

The advantage of dimensional regularization is its translational invariance.
Because of this property, the energy--momentum tensor naively constructed from
bare quantities, Eq.~\eqref{eq:(2.4)}, is conserved and generates
correctly-normalized translations through a WT relation,
\begin{equation}
   \int d^Dx\,\left\langle\partial_\mu T_{\mu\nu}(x)\mathcal{O}\right\rangle
   =-\left\langle\partial_\nu\mathcal{O}\right\rangle,
\label{eq:(2.5)}
\end{equation}
where it is understood that the derivative on the right-hand side is acting all
positions in a gauge-invariant operator~$\mathcal{O}$. Used in combination with
dimensional counting and gauge invariance, this WT relation implies that the
energy--momentum tensor~$T_{\mu\nu}(x)$ is
finite~\cite{Joglekar:1975jm,Nielsen:1977sy} and thus, in the minimal
subtraction (MS) scheme,\footnote{Here, we define the renormalized operator by
subtracting its vacuum expectation value. In the perturbation theory using
dimensional regularization, this subtraction is automatic.}
\begin{equation}
   T_{\mu\nu}(x)-\left\langle T_{\mu\nu}(x)\right\rangle
   =\left\{T_{\mu\nu}\right\}_R(x).
\label{eq:(2.6)}
\end{equation}

The finiteness of the energy--momentum tensor~\eqref{eq:(2.4)} provides further
useful information on the renormalization of dimension~$4$ gauge-invariant
operators. The gauge coupling renormalisation with dimensional regularization
is defined by
\begin{equation}
   g_0^2\equiv\mu^{2\epsilon}g^2Z,
\label{eq:(2.7)}
\end{equation}
where $\mu$ is the renormalization scale and $Z$ is the renormalization factor.
In the MS scheme,
\begin{equation}
   Z=1-\frac{1}{\epsilon}
   \left[b_0g^2+\frac{1}{2}b_1g^4+O(g^6)\right]
   +O\left(\frac{1}{\epsilon^2}\right),
\label{eq:(2.8)}
\end{equation}
and
\begin{equation}
   b_0=\frac{11N}{48\pi^2},\qquad
   b_1=\frac{17N^2}{384\pi^4}.
\label{eq:(2.9)}
\end{equation}
From the rotational invariance that the dimensional regularization keeps, we
see that the operator-renormalization possesses the following
structures:\footnote{Here again, we define renormalized operators by
subtracting their vacuum expectation values.}
\begin{align}
   &F_{\mu\rho}^a(x)F_{\nu\rho}^a(x)
   -\left\langle F_{\mu\rho}^a(x)F_{\nu\rho}^a(x)\right\rangle
\notag\\
   &=Z_T\left\{F_{\mu\rho}^aF_{\nu\rho}^a\right\}_R(x)
   +Z_M\delta_{\mu\nu}\left\{F_{\rho\sigma}^aF_{\rho\sigma}^a\right\}_R(x),
\label{eq:(2.10)}
\end{align}
and
\begin{equation}
   F_{\rho\sigma}^a(x)F_{\rho\sigma}^a(x)
   -\left\langle F_{\rho\sigma}^a(x)F_{\rho\sigma}^a(x)\right\rangle
   =Z_S\left\{F_{\rho\sigma}^aF_{\rho\sigma}^a\right\}_R(x).
\label{eq:(2.11)}
\end{equation}
Substituting the above relations into~Eqs.~\eqref{eq:(2.4)}
and~\eqref{eq:(2.6)}, we have
\begin{align}
   &\left\{T_{\mu\nu}\right\}_R(x)
\notag\\
   &=\frac{1}{g^2}\mu^{-2\epsilon}
   Z^{-1}\biggl[Z_T\left\{F_{\mu\rho}^aF_{\nu\rho}^a\right\}_R(x)
   -\frac{1}{4}(Z_S-4Z_M)
   \delta_{\mu\nu}\left\{F_{\rho\sigma}^aF_{\rho\sigma}^a\right\}_R(x)
   \biggr].
\label{eq:(2.12)}
\end{align}
Since the left-hand side is finite for~$\epsilon\to0$, in the MS scheme in
which only pole terms are subtracted, we infer (by considering the cases,
$\mu\neq\nu$ and~$\mu=\nu$) that
\begin{equation}
   Z_T=Z=1-b_0g^2\frac{1}{\epsilon}+O(g^4)
\label{eq:(2.13)}
\end{equation}
and
\begin{equation}
   Z_S-4Z_M=Z.
\label{eq:(2.14)}
\end{equation}

\subsection{Implications of the trace anomaly}
Another important property of the energy--momentum tensor~\eqref{eq:(2.4)} is
the trace anomaly~\cite{Adler:1976zt,Nielsen:1977sy,Collins:1976yq},
\begin{equation}
   \delta_{\mu\nu}\left\{T_{\mu\nu}\right\}_R(x)
   =-\frac{\beta}{2g^3}
   \left\{F_{\rho\sigma}^aF_{\rho\sigma}^a\right\}_R(x).
\label{eq:(2.15)}
\end{equation}
By Eq.~\eqref{eq:(2.6)}, this relation is equivalent to
\begin{align}
   \delta_{\mu\nu}
   \left[T_{\mu\nu}(x)-\left\langle T_{\mu\nu}(x)\right\rangle\right]
   &=\epsilon\frac{1}{2g_0^2}
   F_{\rho\sigma}^a(x)F_{\rho\sigma}^a(x)
   -\left\langle\epsilon\frac{1}{2g_0^2}
   F_{\rho\sigma}^a(x)F_{\rho\sigma}^a(x)\right\rangle
\notag\\
   &\xrightarrow{\epsilon\to0}
   -\frac{\beta}{2g^3}
   \left\{F_{\rho\sigma}^aF_{\rho\sigma}^a\right\}_R(x).
\label{eq:(2.16)}
\end{align}
In Eqs.~\eqref{eq:(2.15)} and~\eqref{eq:(2.16)}, $\beta$ denotes the
$\beta$~function for $D=4$, defined by
\begin{equation}
   \beta\equiv\left(\mu\frac{\partial}{\partial\mu}\right)_0g
   =-\frac{1}{2}g\left(\mu\frac{\partial}{\partial\mu}\right)_0\ln Z,
\label{eq:(2.17)}
\end{equation}
where the subscript~$0$ implies that the derivative is taken while the bare
quantities are kept fixed. Equations~\eqref{eq:(2.8)} and~\eqref{eq:(2.7)}
yield
\begin{equation}
   \beta=-b_0g^3-b_1g^5+O(g^7).
\label{eq:(2.18)}
\end{equation}
Then, substituting Eqs.~\eqref{eq:(2.12)} into~Eq.~\eqref{eq:(2.15)} and using
Eqs.~\eqref{eq:(2.13)} and~\eqref{eq:(2.14)}, we observe that
\begin{equation}
   \delta_{\rho\lambda}
   \left\{F_{\rho\sigma}^aF_{\lambda\sigma}^a\right\}_R(x)
   =\left(1-\frac{\beta}{2g}\right)
   \left\{F_{\rho\sigma}^aF_{\rho\sigma}^a\right\}_R(x),
\label{eq:(2.19)}
\end{equation}
i.e., the contraction with the metric and the minimal subtraction, the
subtraction of $1/\epsilon$~poles, do not commute; this is a peculiar but
legitimate property of the dimensional regularization~\cite{Collins:1984xc}.

Also, substituting Eqs.~\eqref{eq:(2.7)} and~\eqref{eq:(2.11)}
into~Eq.~\eqref{eq:(2.16)}, we see
\begin{equation}
   \epsilon\frac{Z_S}{Z}
   \xrightarrow{\epsilon\to0}
   -\frac{\beta}{g}.
\label{eq:(2.20)}
\end{equation}
In the MS scheme in which only pole terms are subtracted, this implies
\begin{equation}
   Z_S=\left(1-\frac{\beta}{g}\frac{1}{\epsilon}\right)Z
   =1+O(g^4),
\label{eq:(2.21)}
\end{equation}
and Eq.~\eqref{eq:(2.14)} then shows
\begin{equation}
   Z_M=-\frac{\beta}{4g}Z\frac{1}{\epsilon}
   =\frac{b_0}{4}g^2\frac{1}{\epsilon}+O(g^4).
\label{eq:(2.22)}
\end{equation}
We thus observe that all the renormalization constants
in~Eqs.~\eqref{eq:(2.10)} and~\eqref{eq:(2.11)}, $Z_T$, $Z_M$ and~$Z_S$, in the
MS scheme can eventually be expressed by the gauge coupling renormalization
constant~$Z$ in~Eq.~\eqref{eq:(2.7)}.
